\clearpage\pagestyle{otMain}\setlength{\otBlocEspacement}{2.5mm}\fontsize{11}{15}\selectfont\corrigesChapitre{17}{Solides, disques et volumes}
\begin{blocLivre}{cor-17-1}\begin{corrige}{1}{Reconnaître un solide}\demarche{Observer les bases et les faces latérales.}
A est un $\res{\text{pavé droit}}$ avec six faces rectangulaires. B est un $\res{\text{prisme droit triangulaire}}$ avec deux triangles et trois rectangles, soit cinq faces.
\label{fig-cor-17-1}\figureGeo[corrige]{sol-1}
\end{corrige}\end{blocLivre}
\begin{blocLivre}{cor-17-2}\begin{corrige}{2}{Volume d’un cube}\demarche{Multiplier les trois dimensions.}
$V=4^3=\res{64\ \mathrm{cm}^3}$. On compte quatre cubes unité dans chaque direction.
\label{fig-cor-17-2}\figureGeo[corrige]{sol-2}
\end{corrige}\end{blocLivre}
\begin{blocLivre}{cor-17-3}\begin{corrige}{3}{Volume d’un pavé}\demarche{Utiliser longueur fois largeur fois hauteur.}
$V=7\times5\times3=\res{105\ \mathrm{cm}^3}$. Les trois dimensions sont dans la même unité.
\label{fig-cor-17-3}\figureGeo[corrige]{sol-3}
\end{corrige}\end{blocLivre}
\begin{blocLivre}{cor-17-4}\begin{corrige}{4}{Aire d’un disque}\demarche{Élever le rayon au carré.}
$A=\pi\times4^2=\res{16\pi\ \mathrm{cm}^2}$. La formule utilise le rayon.
\label{fig-cor-17-4}\figureGeo[corrige]{sol-4}
\end{corrige}\end{blocLivre}
\begin{blocLivre}{cor-17-5}\begin{corrige}{5}{Volume et capacité}\demarche{Relier les unités de volume et de capacité.}
2,5 dm$^3=\res{2{,}5\ \mathrm L}$ et $\res{2500\ \mathrm{cm}^3}$. Un dm$^3$ vaut 1000 cm$^3$.
\label{fig-cor-17-5}\figureGeo[corrige]{sol-5}
\end{corrige}\end{blocLivre}
\begin{blocLivre}{cor-17-6}\begin{corrige}{6}{Prisme triangulaire}\demarche{Calculer l’aire de base avant le volume.}
L’aire de la base est $6\times4:2=12$ cm$^2$. Le volume vaut $12\times9=\res{108\ \mathrm{cm}^3}$.
\label{fig-cor-17-6}\figureGeo[corrige]{sol-6}
\end{corrige}\end{blocLivre}
\begin{blocLivre}{cor-17-7}\begin{corrige}{7}{Cylindre}\demarche{Multiplier l’aire du disque par la hauteur.}
$V=\pi\times2^2\times10=\res{40\pi\ \mathrm{cm}^3}$, soit environ $\res{126\ \mathrm{cm}^3}$.
\label{fig-cor-17-7}\figureGeo[corrige]{sol-7}
\end{corrige}\end{blocLivre}
\begin{blocLivre}{cor-17-8}\begin{corrige}{8}{Diamètre donné}\demarche{Passer du diamètre au rayon.}
Le rayon vaut $6:2=3$ cm. Donc $V=\pi\times3^2\times5=\res{45\pi\ \mathrm{cm}^3}$. Utiliser 6 comme rayon multiplierait l’aire par quatre.
\label{fig-cor-17-8}\figureGeo[corrige]{sol-8}
\end{corrige}\end{blocLivre}
\begin{blocLivre}{cor-17-9}\begin{corrige}{9}{Un patron de prisme}\demarche{Faire correspondre les bords des faces.}
Chaque rectangle relie un côté du premier triangle au côté correspondant du second. Leurs dimensions sont $\res{3\times6}$, $\res{4\times6}$ et $\res{5\times6}$ cm. Les deux triangles sont identiques.
\label{fig-cor-17-9}\figureGeo[corrige]{sol-9}
\end{corrige}\end{blocLivre}
\begin{blocLivre}{cor-17-10}\begin{corrige}{10}{Un patron de cylindre}\demarche{Dérouler un tour complet de la base.}
Une dimension est la hauteur : $\res{7\ \mathrm{cm}}$. L’autre est le périmètre de base : $2\pi\times3=\res{6\pi\ \mathrm{cm}}$, environ 18,85 cm. Ajouter deux disques de rayon 3 cm.
\label{fig-cor-17-10}\figureGeo[corrige]{sol-10}
\end{corrige}\end{blocLivre}
\begin{blocLivre}{cor-17-11}\begin{corrige}{11}{Compter des cubes}\demarche{Additionner les hauteurs des piles.}
Chaque nombre donne le nombre de cubes empilés à cette place. Le total vaut $2+1+3+1+2+0=\res9$ cubes, y compris ceux masqués sous les cubes du dessus.
\label{fig-cor-17-11}\figureGeo[corrige]{sol-11}
\end{corrige}\end{blocLivre}
\begin{blocLivre}{cor-17-12}\begin{corrige}{12}{Remplir un aquarium}\demarche{Convertir le volume obtenu en capacité.}
$V=60\times30\times25=45000$ cm$^3$. Comme 1000 cm$^3=1$ L, cela représente $\res{45\ \mathrm L}$. On utilise la hauteur d’eau, pas la hauteur totale du récipient.
\label{fig-cor-17-12}\figureGeo[corrige]{sol-12}
\end{corrige}\end{blocLivre}
\begin{blocLivre}{cor-17-13}\begin{corrige}{13}{Comparer des récipients}\demarche{Calculer chaque capacité avec sa formule.}
$V_A=480$ cm$^3$. $V_B=160\pi\approx502{,}65$ cm$^3$. Le cylindre $\res B$ contient environ $\res{22{,}65\ \mathrm{cm}^3}$ de plus. On compare des volumes dans la même unité.
\label{fig-cor-17-13}\figureGeo[corrige]{sol-13}
\end{corrige}\end{blocLivre}
\begin{blocLivre}{cor-17-14}\begin{corrige}{14}{Compléter une perspective}\demarche{Garder les arêtes fuyantes parallèles et de même longueur.}
Depuis A et C, on trace des arêtes parallèles à [BE] et de même longueur sur le dessin : on obtient D et F, puis la face arrière DEF, identique à ABC. Les arêtes issues de D, sommet caché, sont en pointillés. Les arêtes $\res{\text{[AD] et [CF]}}$ ont la même longueur que [BE] : c’est la longueur du prisme. La face $\res{\text{DEF}}$ est parallèle à ABC ; les faces latérales sont des rectangles.
\label{fig-cor-17-14}\figureGeo[corrige]{sol-14}
\end{corrige}\end{blocLivre}
\begin{blocLivre}{cor-17-15}\begin{corrige}{15}{Une boîte ouverte}\demarche{Séparer le volume intérieur et l’aire du patron.}
La capacité vaut $12\times8\times5=\res{480\ \mathrm{mL}}$. Le carton comprend le fond et quatre côtés : $12\times8+2\times12\times5+2\times8\times5=\res{296\ \mathrm{cm}^2}$. On n’ajoute pas de couvercle.
\label{fig-cor-17-15}\figureGeo[corrige]{sol-15}
\end{corrige}\end{blocLivre}
